\par\Needspace{7\baselineskip}\bigskip
\begin{center}
\pdfbookmark[2]{Gauss to Bessel, Göttingen, 18 May 1814}{pallas-brief-18140518}
\textsc{Gauss} to \textsc{Bessel}, Göttingen, 18 May 1814.
\end{center}
\SourceLeaf{430}\noindent\hspace*{\parindent}. . . . .\quad I have completed the perturbations of Pallas by Jupiter, so far as I wished to calculate them. The first 9 oppositions can already be reconciled fairly well with these perturbations alone, but, strangely enough, \textls[180]{incomparably} better still, so that the errors \textls[180]{average} only \(4''\text{—}5''\), \(\BareRoot{\frac{\text{sum of squared errors}}{\text{number}}}\mathpunct{\raisebox{.45ex}{,}}\) if Jupiter's mass is increased by an appreciable amount. This can be judged better once the perturbations by Saturn, which have also been completed, are taken into account. I have now also begun the perturbations by Mars, which are admittedly numerous but nearly all very small. There is one of \(10''\), however, which in future will provide the best means of determining the mass of Mars.\quad . . . . .
\par
